\section{$4$-dynamic coloring of $C_p(1,3)$}

\begin{theorem}\label{thm:circulant-r4}
Let $p\ge12$. Then
\[
\chi_4^d(C_p(1,3))=
\begin{cases}
5, & p\equiv0\pmod5,\\
8, & p=16,\\
7, & p\in\{13,14,19,21,26,31\},\\
6, & \text{otherwise}.
\end{cases}
\]
\end{theorem}

\begin{proof}
Since $C_p(1,3)$ is $4$-regular, a proper coloring is $4$-dynamic
if and only if vertices at distance at most two receive different
colors. Hence
\[
\chi_4^d(C_p(1,3))=\chi(H_p),
\]
where
$H_p=C_p(1,2,3,4,6)$.
For every $i$, the set
$
\{i,i+1,i+2,i+3,i+4\}
$
induces a $K_5$ in $H_p$. Thus $\chi(H_p)\ge5$. Moreover, in any
$5$-coloring, two consecutive such cliques force
$
\phi(i+5)=\phi(i)
$
for every $i$. Therefore a $5$-coloring can exist only if $5\mid p$.
Consequently,
$
\chi(H_p)\ge6
$
whenever $5\nmid p$.

We need stronger lower bounds only for a few small values of $p$.
Let $S$ be an independent set of $H_p$, and let
$g_1,\ldots,g_t$ be the cyclic distances between consecutive vertices of
$S$. Since the distances $1,2,3,4,6$ are forbidden, every $g_j$ is
either $5$ or at least $7$. Hence
\[
\alpha(H_p)\le\left\lfloor\frac p5\right\rfloor.
\]
Furthermore, if $p=5q+1$, then $\alpha(H_p)\le q-1$, since $q$
distances would require
$
1=\sum_{j=1}^q(g_j-5),
$
where every summand is either $0$ or at least $2$.
Using
$
\chi(H_p)\ge
\left\lceil\frac{p}{\alpha(H_p)}\right\rceil,
$
we obtain
$
\chi(H_{16})\ge8
$
and
$
\chi(H_p)\ge7
\qquad
\text{for }p\in\{13,14,19,21,26,31\}.
$

It remains to give matching colorings. Let
\[
A=01234,\qquad
B=012340521435,\qquad
C=012340153421504235.
\]
If $p=5q$, then $A^q$ is a proper $5$-coloring of $H_p$.

For the nonexceptional values with $5\nmid p$, write $p=5q+r$.
The following cyclic color words use six colors:
\[
\begin{array}{c|c|c}
r & \text{condition} & \phi\\ \hline
1 & q\ge7 & A^{q-7}C^2\\
2 & q\ge2 & A^{q-2}B\\
3 & q\ge3 & A^{q-3}C\\
4 & q\ge4 & A^{q-4}B^2.
\end{array}
\]
A direct verification shows that each word is a proper coloring of
$H_p$.

Finally, let
\[
D=0123456,\qquad E=01234567.
\]
The exceptional cases are covered by
\[
\begin{array}{c|c}
p&\phi\\ \hline
13&D201534\\
14&D^2\\
16&E^2\\
19&D^2A\\
21&D^3\\
26&D^3A\\
31&D^4\,234.
\end{array}
\]
Again, a direct verification shows that these are proper colorings of
$H_p$. They use seven colors, except for $p=16$, where eight colors
are used. Together with the lower bounds above, the result follows.
\end{proof}